\section{Evaluation design}
\label{sec:design}

\subsection{Matched comparison}

The comparison asks whether the 30-second change of a signal at $u$ ranks the run-up to a stall above the run-up to a
pause that did not become one. Controls are the starts $u$ of pauses of at least 2~s and less than 30~s by the same
writer, collected across the writer's assignments. A control is excluded if it falls in $[s - 30\,\mathrm{s}, t_0]$ of
any onset in the same assignment, if condition A is true at its $u$, or if it lies less than 60~s after the start of its
work block; the same 60~s rule applies to stalls. Each stall and control is assigned a \emph{boundary type}, whether the
character at $u$ is a newline (registered as the coarsest position in code, a line break; finer positions are an
exploration in Section~\ref{sec:res-explore}), and a \emph{stratum}: the event stratum if a run, paste or switch occurred in
$[u - 30\,\mathrm{s}, u]$, and the no-event stratum otherwise. Only information up to $u$ is used for the stratum. For
each stall, up to five controls are drawn without replacement from the candidates of the same writer, boundary type and
stratum. The confirmatory test uses the no-event stratum; the event stratum is described. The comparison is thus stratified
within writer, boundary type and stratum rather than matched one to one, and the exclusion of controls near a later
onset uses information after $u$ that a running trigger would not have.

\subsection{Statistic}

For each signal on its own, every pair of a stall and a control drawn within the same writer, boundary type and stratum
scores 1 if the stall's value exceeds the control's by at least $\varepsilon$, 0 if it falls below by at least
$\varepsilon$, and 0.5 otherwise; a control drawn for two stalls of the same cell is counted once. Every stall is
paired with every distinct control drawn in its cell, not only with its own five, so a stall weighs by the number of
controls in its cell (about 156 per stall in the holdout's main stratum). The pooled score over
all pairs of all writers is the \emph{discrimination AUC} (Appendix~\ref{app:definitions} gives the formal definition).
It is 0.5 when the stall's run-up and the control's run-up change by the same amount, even if both rise; differences
that concern only the tails, or that appear only when the two signals are combined, are other estimation targets. Two
tolerances were registered, $\varepsilon = 0.001$ and $\varepsilon = 0$. Writers with more pairs weigh more; we report
the effective number of writers, $(\sum_i n_i)^2 / \sum_i n_i^2$ with $n_i$ the number of pairs of writer $i$, and a
version that weights writers equally. The 95\% interval is the percentile interval of 2,000 resamples of writers, with
each writer's drawn controls held fixed.

Both signals are computed forward in time from the key presses up to $u$, and nothing was fitted to the holdout.
$\theta$ decides when a stall is recognised and thus which pauses count as stalls, and it uses only intervals up to that
time. The intended use is a trigger that fires during a pause, which the AUC alone does not
describe \cite{pelanek2015metrics}. We therefore also report a threshold table: with the 90th percentile of each
writer's drawn controls as the threshold, the share of stalls whose run-up exceeds it and the share of controls that do.

\subsection{Equivalence band and decision}

The band $[0.45, 0.55]$, 0.05 on either side of chance, was set wide after we had seen the matched comparison
in the essay-writing and co-writing data (a wide band makes equivalence easier to declare); it was not derived from a conversion. As a reading aid, an AUC of 0.55
corresponds to a difference in means of about 0.18 standard deviations between two normal distributions of equal
variance. After the band was set, we converted it into the behaviour of a threshold trigger on seen writers
(Appendix~\ref{app:threshold}). The band is wider than the region of practical equivalence of 0.01 used for the
difference in AUC between two models in learning analytics \cite{gardner2018evaluating}; that precedent compares a
different quantity.

Some rules of the main analysis were fixed after a leak had shown that, under an older definition, the AUC on seen
writers lay inside the band (Appendix~\ref{app:prereg}): among them, the stratum window $[u - 30\,\mathrm{s}, u]$,
condition A on both stalls and controls, requiring $\varepsilon = 0$ as well, and the precision rule. With the older
stratum window alone, the holdout gave 0.505 and 0.500 (Section~\ref{sec:res-main}), also inside the band.

For each signal and each $\varepsilon$, the result is \emph{equivalent} if the 95\% interval lies inside the band. This
is the form of the two one-sided tests procedure at one-sided 2.5\% \cite{lakens2017equivalence}. An interval that
crosses 0.55 or 0.45 is undecidable on that side, and an interval entirely above 0.55 or below 0.45 is a difference
beyond the band. An interval wider than 0.10 was registered as insufficient precision, and a signal with more than half
of its pairs tied at $\varepsilon = 0.001$ as degenerate. The test is passed only if all four combinations of two
signals and two tolerances are equivalent.

\subsection{Calibration of the false-equivalence rate}

Because the coverage of resampling by writer is not assured with few or unequal clusters (Section~\ref{sec:related}),
we estimated, on the seen writers, how often the decision declares equivalence when the true AUC lies on the boundary:
provisional signals were shifted until one reached 0.45 or 0.55 while the other stayed at 0.50, and the decision was
applied to 500 draws of 138 writers, the expected size of the holdout. In the version with replacement, which we read
because the holdout is a random sample, the maximum over the configurations examined was 5.2\% (26 of 500; Wilson 95\%
interval 3.6--7.5\%), above the nominal one-sided 2.5\%; without replacement it was at most 1.4\% (Appendix~\ref{app:calibration};
Section~\ref{sec:limitations} gives what the value does not cover). In the holdout every interval ended at
least 0.025 inside the edges of the band, so the decision does not hinge on the exact level.

\subsection{Secondary and descriptive analyses}

A secondary comparison followed the form of earlier reports: every 5~s within a block, a 30~s window ending at $w$ is a
positive if a stall's $u$ falls in the next 60~s, and windows inside long pauses are excluded; the AUC of positives over
negatives within writers is then re-matched by the decile of elapsed time in the block (Appendix~\ref{app:unmatched}).
A descriptive analysis gave 18 standard keystroke quantities, nine levels over the 30~s before $u$ and their changes from
the preceding 30~s, to logistic regression and gradient boosting, trained and evaluated on the same stalls and controls
with folds by writer (Appendix~\ref{app:learners}). Only the learners had a registered positive control: a gate that adds a known difference to one quantity and requires
the learners to recover it (its history is in Appendix~\ref{app:prereg}). The positive control of the estimator was
added afterwards as exploration (Section~\ref{sec:res-bounds}).

\subsection{Status of each analysis}

Each result carries one of the labels in Table~\ref{tab:status}. The confirmation was completed on the holdout in one run, in two
stages recorded on the Open Science Framework, after one interrupted run and one deviation in a verification step
before any value was computed (Appendix~\ref{app:prereg}).

\begin{table}
  \caption{Status of the analyses. Labels are used in the tables of results.}
  \label{tab:status}
  \centering\small
  \begin{tabularx}{\linewidth}{l >{\raggedright\arraybackslash}X}
    \toprule
    Label & Meaning \\
    \midrule
    {[C]} & Confirmatory decision on the CS1 holdout (preregistered rules, deviation, start records and results
    registration). \\
    {[C-D]} & Quantities on the same holdout that the preregistration promised to report but that do not enter the
    decision. \\
    {[D]} & Descriptions on data seen before the rules were frozen. \\
    {[E]} & Exploration outside the registration on seen data, run after the confirmation with conditions committed
    beforehand. \\
    {[E-H]} & Exploration outside the registration on the holdout, after it had been opened, with conditions committed
    beforehand. \\
    \bottomrule
  \end{tabularx}
\end{table}
